\documentclass[12pt]{article}
\usepackage{amsmath, amssymb}
\usepackage[margin=1in]{geometry}
\setlength{\parskip}{6pt}
\title{QUBO Formulation: Network Intrusion Feature Selection}
\date{}
\begin{document}
\maketitle

\subsection*{Network Intrusion Feature Selection}

\paragraph{Problem Statement.}
Given a set of $N$ candidate network-traffic features, the goal is to select a subset of exactly $K$ features for an intrusion detection model. Each feature $i$ has an associated detection relevance score $s_i$, and each pair of features $(i,j)$ has an associated redundancy penalty $p_{ij}$. The objective is to maximize the total detection relevance of the selected features minus the total redundancy among the selected pairs. Formally, we seek a subset of indices of cardinality $K$ that maximizes the net score.

\paragraph{Decision Variables.}
Let $x_i \in \{0,1\}$ be a binary decision variable for each candidate feature $i \in \{0, \dots, N-1\}$, where $x_i = 1$ indicates that feature $i$ is selected, and $x_i = 0$ otherwise.

\paragraph{QUBO Derivation.}
We convert the maximization problem into a minimization problem suitable for a Quadratic Unconstrained Binary Optimization (QUBO) formulation by negating the objective function and incorporating the cardinality constraint as a penalty term.

\textbf{Step 1: Objective Function.}
The original objective is to maximize $S(x) = \sum_{i=0}^{N-1} s_i x_i - \sum_{0 \le i < j \le N-1} p_{ij} x_i x_j$.
Minimizing the negative of this objective yields the unconstrained QUBO component:
\begin{align}
    H_{\text{obj}}(x) &= -\sum_{i=0}^{N-1} s_i x_i + \sum_{0 \le i < j \le N-1} p_{ij} x_i x_j.
\end{align}

\textbf{Step 2: Constraint.}
The requirement to select exactly $K$ features is expressed as:
\begin{align}
    \sum_{i=0}^{N-1} x_i = K.
\end{align}

\textbf{Step 3: Penalty Term.}
We introduce a penalty term with weight $\lambda > 0$ to enforce the constraint. The penalty is proportional to the square of the constraint violation:
\begin{align}
    H_{\text{pen}}(x) &= \lambda \left( \sum_{i=0}^{N-1} x_i - K \right)^2.
\end{align}
Expanding the square and using the idempotent property $x_i^2 = x_i$ for binary variables:
\begin{align}
    \left( \sum_{i=0}^{N-1} x_i - K \right)^2 &= \sum_{i=0}^{N-1} x_i^2 + 2 \sum_{0 \le i < j \le N-1} x_i x_j - 2K \sum_{i=0}^{N-1} x_i + K^2 \nonumber \\
    &= \sum_{i=0}^{N-1} x_i + 2 \sum_{0 \le i < j \le N-1} x_i x_j - 2K \sum_{i=0}^{N-1} x_i + K^2 \nonumber \\
    &= (1 - 2K) \sum_{i=0}^{N-1} x_i + 2 \sum_{0 \le i < j \le N-1} x_i x_j + K^2.
\end{align}
Thus, the penalty contribution is:
\begin{align}
    H_{\text{pen}}(x) &= \lambda (1 - 2K) \sum_{i=0}^{N-1} x_i + 2\lambda \sum_{0 \le i < j \le N-1} x_i x_j + \lambda K^2.
\end{align}

\textbf{Step 4: Combined QUBO Hamiltonian.}
Combining the objective and penalty terms, the total Hamiltonian $H(x) = H_{\text{obj}}(x) + H_{\text{pen}}(x)$ is:
\begin{align}
    H(x) &= \sum_{i=0}^{N-1} \left[ -s_i + \lambda(1 - 2K) \right] x_i + \sum_{0 \le i < j \le N-1} \left[ p_{ij} + 2\lambda \right] x_i x_j + \lambda K^2.
\end{align}

\textbf{Step 5: Q Matrix and Offset.}
The QUBO is defined by $H(x) = \sum_{i,j} Q_{i,j} x_i x_j + c$. Reading off the coefficients from the general form above, we obtain:
\begin{align}
    Q_{i,i} &= -s_i + \lambda(1 - 2K), \quad \forall i \in \{0, \dots, N-1\}, \\
    Q_{i,j} &= \frac{p_{ij}}{2} + \lambda, \quad \forall 0 \le i < j \le N-1, \\
    c &= \lambda K^2.
\end{align}
Here, $Q_{i,j}$ for $i \neq j$ represents the symmetric entry such that the term $x_i x_j$ in the quadratic form $\sum_{i,j} Q_{i,j} x_i x_j$ contributes $2 Q_{i,j} x_i x_j$.

\paragraph{Penalty Weight Justification.}
The penalty weight $\lambda$ must be chosen sufficiently large to ensure that any violation of the cardinality constraint incurs a cost strictly greater than the maximum possible gain in the objective function. Let $\lambda > \max_{i} s_i$. The minimum penalty for a deviation of $\delta = |\sum x_i - K| \ge 1$ is $\lambda \delta^2 \ge \lambda$. The maximum gain in the objective function achievable by adding or removing features is bounded by the sum of relevance scores; however, for a deviation of $\delta$, the gain scales at most linearly with $\delta \cdot \max_i s_i$, while the penalty scales quadratically as $\lambda \delta^2$. Since $\lambda > \max_i s_i$, we have $\lambda \delta^2 > \delta \max_i s_i$ for all $\delta \ge 1$. Thus, any infeasible solution has an energy strictly higher than any feasible solution, ensuring the constraint is satisfied at the optimum.

\paragraph{Correctness Argument.}
Minimizing $H(x)$ over $\{0,1\}^N$ recovers the optimal feature selection. For any feasible solution $x$ (where $\sum x_i = K$), the penalty term vanishes, and $H(x) = -S(x)$. Therefore, among feasible solutions, minimizing $H$ is equivalent to maximizing $S$. For any infeasible solution $x'$, the penalty term is positive. By the choice of $\lambda$, the penalty strictly dominates any potential increase in $S(x')$, implying $H(x') > H(x)$ for any feasible $x$. Consequently, the global minimum of $H(x)$ must occur at a feasible solution that maximizes the original objective.

\paragraph{Illustrative Example.}
Consider an instance with $N=3$ candidate features and $K=2$ features to be selected. The parameters are relevance scores $s_0, s_1, s_2$ and redundancy penalties $p_{01}, p_{02}, p_{12}$. Using the general formulas with $N=3, K=2$, and choosing $\lambda = 1000$ (satisfying $\lambda > \max_i s_i$), we instantiate the QUBO.

The penalty expansion coefficients are:
\begin{align}
    \lambda(1 - 2K) &= 1000(1 - 4) = -3000, \\
    2\lambda &= 2000, \\
    c &= 1000 \cdot 2^2 = 4000.
\end{align}
The resulting Hamiltonian is:
\begin{align}
    H(x) &= (-s_0 - 3000)x_0 + (-s_1 - 3000)x_1 + (-s_2 - 3000)x_2 \nonumber \\
    &\quad + (p_{01} + 2000)x_0 x_1 + (p_{02} + 2000)x_0 x_2 + (p_{12} + 2000)x_1 x_2 + 4000.
\end{align}
The feasible solutions are binary vectors with exactly two ones: $x^{(1)} = (1,1,0)$, $x^{(2)} = (1,0,1)$, and $x^{(3)} = (0,1,1)$. Evaluating $H$ on these yields:
\begin{align}
    H(x^{(1)}) &= -s_0 - s_1 + p_{01}, \\
    H(x^{(2)}) &= -s_0 - s_2 + p_{02}, \\
    H(x^{(3)}) &= -s_1 - s_2 + p_{12}.
\end{align}
The optimal solution $x^*$ is the vector among $\{x^{(1)}, x^{(2)}, x^{(3)}\}$ that minimizes $H(x)$, which corresponds to the subset of two features maximizing $s_i + s_j - p_{ij}$. For instance, if $s = (10, 8, 6)$ and $p = (2, 3, 1)$, then $H(x^{(1)}) = -16$, $H(x^{(2)}) = -13$, $H(x^{(3)}) = -13$, so $x^* = (1,1,0)$ is optimal with energy $-16$.

\end{document}